\chapter{Gauge Fields, Interaction Scales, and Charged Weak Currents}

Quantum electrodynamics and quantum chromodynamics look very different in the
laboratory, yet both are gauge theories. In fact, nearly every interaction in
the Standard Model is organized this way. The word \emph{gauge} has a deeper
meaning involving local symmetry, but we can postpone that machinery. For the
moment we need a minimal physical description: a gauge theory contains
Maxwell-like fields, conserved sources, symmetry transformations, gauge bosons,
and rules for emitting and absorbing those bosons.

\section{Maxwell Theory as the Prototype}

\subsection{Fields and potentials}

The electromagnetic field has six components: three components of electric
field and three components of magnetic field. The same information can be
encoded, with some redundancy, in a four-vector potential
\begin{equation}
A_\mu=(A_0,\boldsymbol A),
\end{equation}
where \(A_0\) is the electrostatic potential and \(\boldsymbol A\) is the
spatial vector potential. In the electrostatic limit the electric field is
obtained from the spatial gradient,
\begin{equation}
\boldsymbol E=-\boldsymbol\nabla A_0,
\qquad
\boldsymbol F=q\boldsymbol E,
\end{equation}
while the magnetic field is
\begin{equation}
\boldsymbol B=\boldsymbol\nabla\times\boldsymbol A.
\end{equation}
The sign in the first equation depends on the standard definition of the
potential; the board discussion is concerned mainly with the fact that fields
are derivatives of potentials.
\editorialnote{For time-dependent fields the complete relation is
\(\boldsymbol E=-\boldsymbol\nabla A_0-\partial_t\boldsymbol A\). The lecture
deliberately works at the simpler electrostatic level here.}

\begin{figure}[htbp]
\centering
\includegraphics[width=0.62\textwidth]{lecture_05_figure_02.png}
\caption{The electromagnetic four-potential and its derivative relation to the
field.}
\lectureframe{00:03:33}
\end{figure}

Thus electrodynamics admits two equivalent descriptions: one in terms of
\(\boldsymbol E\) and \(\boldsymbol B\), and one in terms of \(A_\mu\).
The latter contains only four displayed components because different
potentials can describe the same physical fields. That redundancy is the first
hint behind the word gauge, though we will not use it formally yet.

\subsection{Waves and polarization}

Maxwell's equations support waves. Take a plane wave moving along a unit vector
\(\hat{\boldsymbol k}\). At every point the electric and magnetic fields are
transverse:
\begin{equation}
\boldsymbol E\cdot\hat{\boldsymbol k}=0,
\qquad
\boldsymbol B\cdot\hat{\boldsymbol k}=0,
\qquad
\boldsymbol E\cdot\boldsymbol B=0.
\end{equation}
In units with \(c=1\), their magnitudes are equal,
\begin{equation}
|\boldsymbol E|=|\boldsymbol B|.
\end{equation}
For a linearly polarized wave, the direction of \(\boldsymbol E\) specifies
the polarization. The entire pattern propagates at the speed of light.

Quantization changes the language but not the polarization structure. A
quantum of the electromagnetic wave is a photon. Drawing a photon as a point
is inevitably misleading, but it is useful to imagine that the point carries
a small transverse pointer recording its polarization. A quarter-wave plate
or another optical element can rotate that polarization. A static Coulomb
field may likewise be pictured, only roughly, as arising from repeated
emission and absorption of photons.

\subsection{Sources, flux, and conservation}

Electromagnetic waves can propagate through empty space, but the field also
has sources: electric charges and currents. A positive point charge produces
an outward radial electric field proportional to \(1/r^2\). The number of
field lines piercing any sphere around the charge is independent of the
sphere's radius. Mathematically,
\begin{equation}
\oint_{S^2}\boldsymbol E\cdot d\boldsymbol S
=\frac{Q_{\mathrm{enc}}}{\epsilon_0},
\label{eq:l5-gauss}
\end{equation}
in SI units.

This is more than a convenient way to calculate the field. Electric field
lines may end only on charge. If one isolated charge simply disappeared, its
field would either have to disappear everywhere at once, violating causality,
or a front of missing field would have to propagate outward with field lines
ending on empty space. Neither is allowed. The geometry of flux therefore
makes charge conservation natural: charge can move, and opposite charges can
annihilate, but net charge cannot vanish locally without a compensating
current.

The corresponding local statement is the continuity equation
\begin{equation}
\partial_t\rho+\boldsymbol\nabla\cdot\boldsymbol J=0.
\label{eq:l5-continuity}
\end{equation}
\editorialnote{Equation~\eqref{eq:l5-continuity} is the standard local
mathematical form of the causal flux argument developed verbally in the
lecture.}

\subsection{A working definition of gauge theory}

We can now state the operational pattern. A gauge theory has:
\begin{itemize}
\item electric-like and magnetic-like fields;
\item waves that, when interactions are weak enough, resemble light waves;
\item sources analogous to electric charge;
\item a Gauss law tying flux to those sources;
\item a conservation law and an associated symmetry;
\item quanta of the field, called gauge bosons; and
\item a coupling constant controlling emission and absorption.
\end{itemize}
The source need not be electric charge. As a deliberately unphysical example,
imagine gauging baryon number. Protons and neutrons both have baryon number
\(+1\), so they would repel through a new Coulomb-like force; a proton and an
antiproton would attract. Such a theory is mathematically imaginable, but
nature supplies no observed long-range gauge field coupled to baryon number.
Conservation by itself does not guarantee a gauge force.

Electric charge does come with a symmetry. For a charged matter field,
\begin{equation}
\psi(x)\longmapsto e^{i\theta}\psi(x).
\label{eq:l5-u1}
\end{equation}
The phase symmetry and charge conservation are two sides of the same
structure. Later we can make \(\theta\) depend on spacetime and see why the
potential \(A_\mu\) is required. Here the global transformation is enough to
identify the conserved charge.

\section{QCD as the Non-Abelian Example}

\subsection{Triplets, anti-triplets, and the gluon matrix}

In quantum chromodynamics a quark field carries a color index
\(i\in\{r,g,b\}\) and transforms in the fundamental representation of
\(SU(3)\):
\begin{equation}
q'_i=U_{ij}q_j,
\qquad
U^\dagger U=\boldsymbol 1,
\qquad
\det U=1.
\label{eq:l5-q-transform}
\end{equation}
This is a rotation in a three-dimensional complex color space, not a rotation
of ordinary position space. The three-component representation is also called
the defining representation.

Antiparticle fields are represented by the complex-conjugate variables. Taking
the complex conjugate of Equation~\eqref{eq:l5-q-transform} gives
\begin{equation}
q_i^{\prime *}=U_{ij}^*q_j^*.
\label{eq:l5-antiquark-transform}
\end{equation}
Thus quarks transform as \(\boldsymbol 3\), whereas antiquarks transform as
\(\bar{\boldsymbol 3}\).

\begin{figure}[htbp]
\centering
\includegraphics[width=0.62\textwidth]{lecture_05_figure_03.png}
\caption{The fundamental transformation and its complex-conjugate
anti-triplet transformation.}
\lectureframe{00:18:47}
\end{figure}

\begin{classroomqa}
\audiencequestion{How many \(U\) matrices are there -- seven, eight, or nine?
\lecturetimestamp{00:19:37}}
\lecturerresponse{There are infinitely many group elements \(U\). The finite
number being counted is the number of independent generators, which is eight
for \(SU(3)\). A general \(3\times3\) matrix has nine matrix positions, but
special unitarity leaves eight continuous directions.}
\end{classroomqa}

The gluon field may be written as a traceless matrix \(A_{ij}\). One index
transforms like a quark and the other like an antiquark, so its color structure
is that of
\begin{equation}
\boldsymbol 3\otimes\bar{\boldsymbol 3}
=\boldsymbol 1\oplus\boldsymbol 8.
\end{equation}
QCD retains the traceless octet,
\begin{equation}
A_{ii}=\operatorname{tr}A=0,
\end{equation}
and discards the singlet. That is why QCD has eight gluons rather than nine.
More generally, \(SU(N)\) has
\begin{equation}
N^2-1
\end{equation}
generators and the same number of gauge-boson fields.

The photon is neutral under electric \(U(1)\): rotating the phase of an
electron does not rotate the photon into a different photon. Gluons are
different. A gluon carries a color and an anticolor, so it transforms in the
adjoint representation. A red--antiblue combination, for example, is not a
color singlet even though it contains one color and one anticolor.

\subsection{Index flow and self-interaction}

The interaction rule can be read by following indices. Schematically, a quark
of color \(i\) can become a quark of color \(j\) while emitting a gluon carrying
the corresponding color flow:
\begin{equation}
q_i\longrightarrow q_j+A_{ij}.
\label{eq:l5-quark-gluon-vertex}
\end{equation}
The precise placement of \(i\) and \(j\) depends on the arrow convention; the
invariant content is that every index entering a vertex must continue through
another line. No color index is lost.

This rule applies in any \(SU(N)\) gauge theory. Fundamental matter has one
index, antifundamental matter has one conjugate index, and a gauge boson has
one of each. The physical names can change, but the index structure remains.

Because a gluon itself carries color, it can emit another gluon. An
\(A_{ij}\) gluon can split through an intermediate index \(k\), schematically,
\begin{equation}
A_{ij}\longrightarrow A_{ik}+A_{kj}.
\label{eq:l5-three-gluon-index}
\end{equation}
This is not a complete Yang--Mills vertex, but it displays the color flow.
Nothing analogous occurs for photons in ordinary electrodynamics. Gluons
therefore exert forces on gluons, and non-Abelian field dynamics is much more
complicated than Maxwell's linear theory.

One qualitative consequence is confinement. Color flux lines cannot end, and
their self-interaction draws them into a tube joining a quark to an antiquark.
The tube carries an approximately constant energy per unit length:
\begin{equation}
V(R)\simeq \sigma R.
\end{equation}
Pulling the endpoints farther apart therefore stores more energy instead of
liberating isolated color. This is the gooey piece of glue picture: the
field itself prevents colored constituents from escaping as free particles.

\section{Couplings and the Meaning of Strong or Weak}

\subsection{Charge as an emission amplitude}

A gauge theory still needs a coupling constant. In QED the electron charge
\(e\) is the coefficient associated with an electron--photon vertex. It is an
amplitude, and probabilities are squares of amplitudes. Operationally,
imagine stopping a fast electron abruptly in a cathode-ray tube. The
probability for radiating a photon is proportional to \(e^2\), with conventional
factors absorbed into the fine-structure constant,
\begin{equation}
\alpha=\frac{e^2}{4\pi\hbar c}
\simeq \frac{1}{137}.
\label{eq:l5-alpha}
\end{equation}
In natural units \(\hbar=c=1\), the coupling is dimensionless. A one-percent
scale for an elementary emission makes electromagnetism weak in the
perturbative sense: most stopped electrons do not each release a spray of
hard photons.

\subsection{Natural units connect energy, time, and distance}

Set
\begin{equation}
c=1,
\qquad
\hbar=1.
\end{equation}
Then time and distance have the same units, while energy has inverse units:
\begin{equation}
E=\hbar\omega\longrightarrow E=\omega\sim\frac1t,
\qquad
x\sim t\sim\frac1E.
\end{equation}
Large energies probe short distances and short times.

For orientation,
\begin{equation}
1\ {\rm eV}^{-1}
=\frac{\hbar c}{1\ {\rm eV}}
\simeq 1.97\times10^{-7}\ {\rm m}.
\end{equation}
An atomic length of order \(10^{-10}\ {\rm m}\) therefore corresponds roughly
to an inverse kiloelectronvolt. These are order-of-magnitude conversions, not
precision spectroscopy.

\subsection{The hierarchy inside an atom}

For a typical atom, take
\begin{equation}
R_{\rm atom}\sim10^{-10}\ {\rm m}.
\end{equation}
Light crosses this distance in
\begin{equation}
t_{\rm transit}^{\rm atom}\sim10^{-18}\ {\rm s}.
\end{equation}
An electron, however, moves at only a small fraction of the speed of light in
a light atom. In the Bohr picture \(v/c\sim\alpha\), so a typical orbital time
is longer:
\begin{equation}
t_{\rm orbit}^{\rm atom}\sim10^{-16}\ {\rm s}.
\end{equation}

An excited atom takes longer still to radiate:
\begin{equation}
t_{\rm decay}^{\rm atom}\sim10^{-9}\ {\rm s}
\end{equation}
as a broad lecture-scale estimate. There are two suppressions in the physical
picture. The electromagnetic force produces only modest acceleration, and an
accelerated electron has another small emission probability controlled by
\(\alpha\). The small fine-structure constant therefore creates a hierarchy
among transit, orbital, and radiative-decay times.

\subsection{The absence of hierarchy in a hadron}

A hadron is an atom-like bound system made of quarks and gluons, but its scale
is much smaller:
\begin{equation}
R_{\rm hadron}\sim10^{-15}\ {\rm m},
\qquad
t_{\rm transit}^{\rm hadron}\sim10^{-23}\ {\rm s}.
\end{equation}
Quarks cross the hadron on essentially the same time scale, and strongly
decaying excited hadrons also live for roughly
\begin{equation}
t_{\rm orbit}^{\rm hadron}
\sim t_{\rm decay}^{\rm hadron}
\sim10^{-23}\ {\rm s}.
\end{equation}
A struck hadron may vibrate and emit a pion almost as quickly as causality
allows. There is no seven- or nine-order-of-magnitude separation like the one
inside an atom.

\begin{classroomqa}
\audiencequestion{Does special relativity imply that a hadron cannot decay
substantially faster than its light-crossing time?
\lecturetimestamp{00:47:15}}
\lecturerresponse{That is a fair way to state the order-of-magnitude bound.
Strong decays occur about as fast as the size of the system permits; there is
no small coupling or other large suppression slowing them down.}
\end{classroomqa}

The corresponding strong coupling is therefore much larger than
\(\alpha\). At an illustrative hadronic scale the lecture quotes
\begin{equation}
\alpha_{\rm QCD}\sim0.2,
\end{equation}
much closer to unity than \(1/137\). This is the dynamical origin of the older
name \emph{strong interactions}.
\editorialnote{The QCD coupling runs with energy, so \(\alpha_s\) is not one
fixed number. The quoted value is an illustrative scale-dependent value, as
the lecture intends.}

\begin{classroomqa}
\audiencequestion{Is the strong analogue simply called the strong-force
constant?
\lecturetimestamp{00:48:17}}
\lecturerresponse{It is conventionally called \(\alpha_{\rm QCD}\), or more
often \(\alpha_s\). The electromagnetic one is usually just \(\alpha\), though
\(\alpha_{\rm QED}\) can be used when both appear in the same discussion.}
\end{classroomqa}

\begin{classroomqa}
\audiencequestion{Would a paper distinguish a bare electromagnetic coupling
from the measured, shielded value?
\lecturetimestamp{00:48:38}}
\lecturerresponse{Yes. One speaks of a bare coupling and a renormalized
coupling. Vacuum polarization makes the effective coupling depend on the
scale at which it is measured.}
\end{classroomqa}

\section{The Puzzle of Weak Decay}

Weak interactions were named from their very long decay times. They are
usually slower than strong decays and can even be slower than atomic
radiation. The oldest example is nuclear beta decay.

\subsection{Neutron beta decay}

A free neutron decays through
\begin{equation}
n\longrightarrow p+e^-+\bar\nu_e.
\label{eq:l5-neutron-decay-intro}
\end{equation}
The lecture uses a lifetime of roughly twelve minutes, absurdly long compared
with \(10^{-23}\ {\rm s}\). Modern precision values do not matter for the
argument.

Part of the slowness is kinematic. Roughly,
\begin{equation}
\begin{aligned}
m_n&\simeq940\ {\rm MeV},
&
m_p&\simeq939\ {\rm MeV},
\\
m_e&\simeq0.5\ {\rm MeV}.
\end{aligned}
\end{equation}
Only a tiny fraction of the neutron's rest energy is available as kinetic
energy. If the neutron were lighter than the sum of the final rest masses,
decay would be impossible. If it lay only infinitesimally above threshold, the
available phase space would be tiny and the decay would be slow. But phase
space alone cannot explain a lifetime measured in minutes.

\subsection{The charged pion removes the threshold excuse}

Consider also
\begin{equation}
\pi^-\longrightarrow e^-+\bar\nu_e
\end{equation}
and its charge-conjugate process. The pion has ample energy above the electron
channel, yet its weak lifetime is still of nanosecond order rather than
hadronic \(10^{-23}\)-second order. A similar channel with a muon will soon
become important. The puzzle is therefore intrinsic to the weak interaction,
not merely an accident of the neutron's near-degenerate masses.

\editorialnote{The charged-pion lifetime is about \(26\ {\rm ns}\), and the
dominant channel is \(\pi^-\to\mu^-\bar\nu_\mu\). The electron channel is
helicity suppressed. The lecture first uses the electron channel because its
particle bookkeeping is simplest, then returns to the dominant muon channel.}

The term beta decay originates historically from calling emitted electrons
beta rays. The pion and neutron examples involve the same charged-current
interaction even though their kinematics differ sharply.

\section{A Table of Quarks and Leptons}

\subsection{Color acts vertically}

Arrange the six quark flavors in three paired columns,
\begin{equation}
(u,d),\qquad(c,s),\qquad(t,b),
\end{equation}
and place a red, green, and blue copy of every quark in separate rows. Color
\(SU(3)\) acts vertically: it rotates red, green, and blue versions of an up
quark into one another, and performs the same color rotation simultaneously
on down, charm, strange, top, and bottom. It does not turn an up quark into a
down quark.

\subsection{Weak isospin acts horizontally}

The weak symmetry acts in the other direction. It mixes each neighboring pair
\begin{equation}
Q_1=\begin{pmatrix}u\\d\end{pmatrix},
\qquad
Q_2=\begin{pmatrix}c\\s\end{pmatrix},
\qquad
Q_3=\begin{pmatrix}t\\b\end{pmatrix}.
\label{eq:l5-quark-doublets}
\end{equation}
The up-type quarks have charge \(+2/3\), and the down-type quarks have charge
\(-1/3\), so the charge changes by one unit across each pair.
\editorialnote{The source momentarily says that a down quark has charge
\(-2/3\). The surrounding subtraction and every subsequent process require
the standard value \(-1/3\), used here.}

Because the symmetry acts on doublets, the natural group is \(SU(2)\). It acts
on every color copy of each quark doublet without changing color.

There is also a fourth row of particles that do not carry color and therefore
do not interact with gluons: the leptons. They form three more doublets,
\begin{equation}
L_e=\begin{pmatrix}\nu_e\\e^-\end{pmatrix},
\qquad
L_\mu=\begin{pmatrix}\nu_\mu\\\mu^-\end{pmatrix},
\qquad
L_\tau=\begin{pmatrix}\nu_\tau\\\tau^-\end{pmatrix}.
\label{eq:l5-lepton-doublets}
\end{equation}
The neutrinos are electrically neutral; the electron, muon, and tau have
charge \(-1\). Again the two members differ by one unit of electric charge.
The electron, muon, and tau have the same gauge quantum numbers but different
masses, just as the three generations of up-type or down-type quarks repeat
the same charge pattern.

\begin{figure}[htbp]
\centering
\includegraphics[width=0.82\textwidth]{lecture_05_figure_05.png}
\caption{The board table: color acts among the red, blue, and green rows,
whereas weak \(SU(2)\) links the paired flavor columns and the lepton row.}
\lectureframe{01:03:35}
\end{figure}

Calling leptons a fourth color is only an analogy at this stage. They fill
a fourth row in the weak-doublet table, but ordinary QCD \(SU(3)\) does not
rotate quarks into leptons.

\section{Constructing the Charged Weak Bosons}

\subsection{Quantum numbers of the charged bosons}

Gauge bosons carry the transformation properties needed to move from one
member of a doublet to the other. Begin with the lepton pair
\(e^-\bar\nu_e\). It has electric charge \(-1\), so the corresponding charged
gauge boson is called \(W^-\):
\begin{equation}
W^-\sim e^-\bar\nu_e.
\end{equation}
The same quantum numbers arise from a down quark and an anti-up quark:
\begin{equation}
W^-\sim d\bar u.
\end{equation}
The symbol \(\sim\) means has the same relevant gauge quantum numbers as.
It does not mean that the \(W\) is literally a bound state of these particles.
The antiparticle has the conjugate combinations,
\begin{equation}
W^+\sim e^+\nu_e\sim\bar d u.
\end{equation}

The charged weak vertices include
\begin{align}
e^-&\longrightarrow\nu_e+W^-,
&
\mu^-&\longrightarrow\nu_\mu+W^-,
\label{eq:l5-lepton-w}\\
d&\longrightarrow u+W^-,
&
s&\longrightarrow c+W^-,
\label{eq:l5-quark-w}
\end{align}
with analogous transitions for the third generation when kinematically
allowed. The gauge boson performs dynamically the same horizontal mixing that
the symmetry performs on the doublet.

\begin{classroomqa}
\audiencequestion{Charge conservation fixes the emitted boson to have charge
\(-1\), but why must the electron, down-quark, and strange-quark transitions
all involve the same particle?
\lecturetimestamp{01:12:56}}
\lecturerresponse{That is an empirical fact encoded by the common weak
symmetry. Nature could consistently have assigned a different \(W\)-type
boson to every family or to quarks and leptons separately. Instead, one
\(SU(2)\) acts simultaneously on all the doublets, and the same charged gauge
bosons mediate all these transitions.}
\end{classroomqa}

\subsection{Reading and reversing a vertex}

A single vertex represents a family of processes. Reversing an incoming line
to an outgoing line changes a particle into its antiparticle, and reversing a
\(W^-\) line gives a \(W^+\). One can therefore read a vertex as emission,
absorption, decay, annihilation, or the inverse reaction, provided charge,
momentum, and all relevant quantum numbers are followed consistently.

\begin{classroomqa}
\audiencequestion{Can the arrows and lines in these weak vertices be flipped?
\lecturetimestamp{01:14:51}}
\lecturerresponse{Yes. A line may be crossed from the final state to the
initial state or vice versa, but a particle then becomes its antiparticle.
The same local vertex consequently describes many differently oriented
reactions.}
\end{classroomqa}

\begin{classroomqa}
\audiencequestion{Does writing a \(W^-\) with strange and charm lines mean
that the boson is literally a strange quark bound to an anti-charm quark?
\lecturetimestamp{01:15:43}}
\lecturerresponse{No. The line bookkeeping says only that the \(W^-\) carries
the quantum-number change between those columns. If the diagram is read as a
decay, a sufficiently energetic \(W^-\) can produce a strange quark and an
anti-charm quark; it is not composed of them.}
\end{classroomqa}

\begin{classroomqa}
\audiencequestion{Can an electron and an antineutrino make a \(W\), which then
produces a quark pair?
\lecturetimestamp{01:17:22}}
\lecturerresponse{Exactly. If enough center-of-mass energy is available, the
initial leptons can form the charged-current intermediate state and that state
can emerge through an allowed quark channel. The two vertices are the same
ones used in decay, read in another orientation.}
\end{classroomqa}

\subsection{The missing third boson}

An \(SU(2)\) group has three generators, hence three gauge fields. So far we
have displayed only \(W^+\) and \(W^-\). A third, electrically neutral field
must be associated with diagonal combinations such as electron--positron and
neutrino--antineutrino quantum numbers. It is deliberately left lurking in the
background. Its relation to the physical \(Z\) boson and photon requires
electroweak mixing, which comes later.

\begin{classroomqa}
\audiencequestion{Were the weak and strong forces suspected at different
historical times?
\lecturetimestamp{01:24:11}}
\lecturerresponse{Their experimental roots appeared together in early
radioactivity. Alpha decay revealed nuclear and strong-interaction physics,
beta decay revealed the weak interaction, and gamma decay emitted
electromagnetic photons. Becquerel's era already contained all three clues,
long before their modern field theories existed.}
\end{classroomqa}

\section{Worked Charged-Current Decays}

\subsection{Pion decay}

The negatively charged pion has valence-quark content
\begin{equation}
\pi^-=d\bar u.
\end{equation}
Those quantum numbers can annihilate into a \(W^-\), which can then produce a
charged lepton and the corresponding antineutrino:
\begin{align}
\pi^-&\longrightarrow W^{-*}
\longrightarrow e^-+\bar\nu_e,
\label{eq:l5-pion-electron}\\
\pi^-&\longrightarrow W^{-*}
\longrightarrow \mu^-+\bar\nu_\mu.
\label{eq:l5-pion-muon}
\end{align}
The star reminds us that the intermediate \(W\) is virtual. The second channel
is the dominant one.

\begin{figure}[htbp]
\centering
\includegraphics[width=0.72\textwidth]{lecture_05_figure_04.png}
\caption{A \(d\bar u\) pion annihilates into a virtual \(W^-\), followed by a
charged lepton and its matching antineutrino.}
\lectureframe{01:19:38}
\end{figure}

\begin{classroomqa}
\audiencequestion{Is this really a Feynman diagram even though the time
direction was drawn horizontally rather than vertically?
\lecturetimestamp{01:20:21}}
\lecturerresponse{Yes. The orientation of time on the page is a convention.
Once one direction is declared to be time and all arrows are interpreted
consistently, the drawing carries the same interaction information.}
\end{classroomqa}

\begin{classroomqa}
\audiencequestion{Could the diagram be read more like a reversible chemical
reaction than as an irreversible one-way decay?
\lecturetimestamp{01:21:35}}
\lecturerresponse{The inverse process is allowed. A pion produced in empty
space decays because the outgoing particles separate and do not ordinarily
recombine. But if a muon and antineutrino are fired together with the right
energy and momentum, there is a nonzero amplitude to form the pion through the
reverse interaction.}
\end{classroomqa}

\subsection{Energy and a virtual intermediate state}

For a pion at rest,
\begin{equation}
m_\pi\simeq140\ {\rm MeV},
\qquad
m_\mu\simeq106\ {\rm MeV},
\qquad
m_\nu\approx0
\end{equation}
at the precision relevant here. The remaining energy becomes kinetic energy
of the outgoing particles. Momentum conservation requires
\begin{equation}
\boldsymbol p_\mu+\boldsymbol p_{\bar\nu}=0,
\end{equation}
so the two particles leave back to back in the pion rest frame.

\begin{classroomqa}
\audiencequestion{How can energy be conserved when the intermediate \(W\) is
far heavier than the pion?
\lecturetimestamp{01:25:08}}
\lecturerresponse{Only the initial and final particles must be on their mass
shell. Four-momentum is conserved at every vertex, but an internal virtual
line need not obey \(E^2=\boldsymbol p^2+m_W^2\). The larger its off-shell
mismatch, the more strongly its propagator suppresses the process.\editorialnote{The lecture uses the traditional energy--time uncertainty
language and says the intermediate state may violate energy briefly. The
modern precise statement is the one in the response: energy-momentum is
conserved, while a virtual internal line need not satisfy the real-particle
mass-shell condition.}}
\end{classroomqa}

\subsection{Muon decay}

A muon can first change into its own neutrino by emitting a virtual \(W^-\):
\begin{equation}
\mu^-\longrightarrow\nu_\mu+W^{-*}.
\end{equation}
The virtual boson then produces an electron and electron antineutrino:
\begin{equation}
\mu^-\longrightarrow
\nu_\mu+e^-+\bar\nu_e.
\label{eq:l5-muon-decay}
\end{equation}
This is a purely leptonic process. The muon is about two hundred times heavier
than the electron, so the final rest masses leave ample kinetic energy.

\begin{figure}[htbp]
\centering
\includegraphics[width=0.72\textwidth]{lecture_05_figure_06.png}
\caption{Muon decay through a virtual charged weak boson.}
\lectureframe{01:29:24}
\end{figure}

\begin{classroomqa}
\audiencequestion{Is the neutral particle emitted with the electron an
antineutrino?
\lecturetimestamp{01:29:06}}
\lecturerresponse{Yes. The full final state contains one muon neutrino, one
electron antineutrino, and the electron. Keeping the flavor labels prevents
the two neutral particles from being confused.}
\end{classroomqa}

The electron cannot reverse this decay in empty space because it is lighter
than the muon. The tau is heavier and can decay through analogous weak
channels. Stability here is partly an accident of the mass ordering, not a
difference in the underlying vertex.

\begin{classroomqa}
\audiencequestion{Is a linear accelerator effectively increasing the
electron's mass so that otherwise forbidden reactions can run in the opposite
direction?
\lecturetimestamp{01:30:23}}
\lecturerresponse{It does not change the electron's invariant rest mass, but
it supplies kinetic energy. Sufficient center-of-mass energy can create heavier
final particles, so the proposed physical idea is correct once mass is
replaced by available collision energy.}
\end{classroomqa}

\subsection{Neutron beta decay at quark level}

The neutron and proton valence compositions are
\begin{equation}
n=udd,
\qquad
p=uud.
\end{equation}
One down quark in the neutron makes the charged-current transition
\begin{equation}
d\longrightarrow u+W^{-*},
\end{equation}
turning \(udd\) into \(uud\). The virtual boson then decays:
\begin{equation}
W^{-*}\longrightarrow e^-+\bar\nu_e.
\end{equation}
Combining the two vertices reproduces
\begin{equation}
n\longrightarrow p+e^-+\bar\nu_e.
\end{equation}

The opposite quark transition \(u\to d+W^+\) would leave a three-down-quark
baryon, but that available state is too heavy for a free neutron to produce.
The actual down-to-up route is the energetically allowed one.

\begin{classroomqa}
\audiencequestion{Could a neutron decay into a proton, a muon, and the
corresponding antineutrino instead?
\lecturetimestamp{01:32:38}}
\lecturerresponse{Not for a free neutron at rest: the neutron--proton mass
difference is far smaller than the muon mass. With enough additional collision
or excitation energy, the same weak vertex can participate in a channel that
produces a muon. The prohibition is kinematic, not a missing interaction.}
\end{classroomqa}

\section{The Unresolved Suppression}

The charged weak vertices explain which processes can occur, but not yet why
they are so slow. One tempting answer would be an extraordinarily tiny weak
coupling. That answer is wrong: the dimensionless weak gauge coupling is of
roughly electromagnetic size. The suppression must come from something else.

The missing ingredient is already visible in the virtual propagator: the
\(W\) is very massive. At energies far below \(m_W\), charged-current
amplitudes are suppressed roughly by powers of \(1/m_W^2\). That connection
will be developed after the Higgs mechanism and electroweak symmetry breaking
explain how the gauge bosons acquire mass.
\editorialnote{The lecture preserves this as the puzzle for the next meeting.
The heavy-\(W\) explanation is included here as a forward connection, not as a
claim that its derivation occurred in this lecture.}

\begin{classroomqa}
\audiencequestion{When we say that the theory has this symmetry, do we mean
that its Lagrangian is invariant under the corresponding operations?
\lecturetimestamp{01:34:14}}
\lecturerresponse{Exactly. The statement that the theory has the symmetry is
shorthand for invariance of the Lagrangian, or equivalently the action, under
those transformations. Writing the full Lagrangian is being postponed, not
replaced by a different definition.}
\end{classroomqa}

The route is now complete: Maxwell theory supplies the prototype; QCD shows
how non-Abelian gauge bosons carry and exchange charge; coupling constants
explain the contrast between atomic and hadronic time scales; and weak
\(SU(2)\) organizes quarks and leptons into doublets connected by
\(W^\pm\). The remaining task is to understand why these gauge-mediated
processes are weak in practice even though their underlying coupling is not
exceptionally small.
